\documentclass[12pt]{article}

% -------------------- Idioma y codificación --------------------
\usepackage[spanish]{babel}
\addto\captionsspanish{\renewcommand{\tablename}{Tabla}}
\usepackage[utf8]{inputenc}
\usepackage[T1]{fontenc}

% -------------------- Formato de página --------------------
\usepackage{geometry}
\geometry{margin=2.5cm}
\usepackage{setspace}
\setstretch{1.15}

% -------------------- Colores --------------------
\usepackage{xcolor}
\definecolor{tituloazul}{RGB}{31,100,160}
\definecolor{barraazul}{RGB}{210,230,250}

% -------------------- Links clickeables --------------------
\usepackage[
    colorlinks=true,
    linkcolor=tituloazul,
    citecolor=tituloazul,
    urlcolor=tituloazul
]{hyperref}

% -------------------- Estilo de secciones --------------------
\usepackage{titlesec}

\titleformat{\section}
{\color{tituloazul}\large\bfseries}
{\thesection}{1em}{}

\titleformat{\subsection}
{\color{tituloazul}\normalsize\bfseries}
{\thesubsection}{1em}{}

\titleformat{\subsubsection}
{\color{tituloazul}\normalsize\bfseries}
{\thesubsubsection}{1em}{}

% -------------------- Encabezados automáticos --------------------
\usepackage{fancyhdr}
\pagestyle{fancy}

\fancyhf{}
\fancyhead[L]{\textcolor{tituloazul}{Electrónica de Potencia}}
\fancyhead[R]{\textcolor{tituloazul}{Práctica 2}}
\fancyfoot[C]{\thepage}

\setlength{\headheight}{15pt}

% -------------------- Figuras --------------------
\usepackage{graphicx}
\graphicspath{{figuras/}}

\usepackage{caption}
\usepackage{subcaption}
\usepackage{float}

% -------------------- Matemáticas --------------------
\usepackage{amsmath}

\usepackage[
    spanish,
    capitalize,
    nameinlink
]{cleveref}

\crefname{figure}{figura}{figuras}
\Crefname{figure}{Figura}{Figuras}

\crefname{equation}{ecuación}{ecuaciones}
\Crefname{equation}{Ecuación}{Ecuaciones}

\crefname{table}{tabla}{tablas}
\Crefname{table}{Tabla}{Tablas}

% -------------------- Unidades --------------------
\usepackage{siunitx}

\sisetup{
    per-mode=symbol,
    detect-all
}

% -------------------- Tablas profesionales --------------------
\usepackage{booktabs}

% -------------------- Referencias IEEE --------------------
\usepackage{cite}

% -------------------- Circuitos --------------------
\usepackage[siunitx]{circuitikz}
\usetikzlibrary{babel}

% -------------------- Símbolos --------------------
\usepackage{latexsym}

% -------------------- Autores --------------------
\newcommand{\estudiantes}{Jose Leonardo Perez Poveda, Sergio Hoyos Mejía, Alberto Salazar}
% ============================================================
%                         DOCUMENTO
% ============================================================

\begin{document}

% -------------------- Título --------------------
\begin{center}

    {\Huge
    \textcolor{tituloazul}{
    \textbf{Práctica 2}\\
    Interruptores de Electrónica de Potencia y
    Circuitos de Disparo. Parte 2
    }}

    \vspace{0.6cm}

    {\textbf{\estudiantes}}

    \vspace{0.5cm}

    \textit{Ingeniería Electrónica}\\
    \textit{Universidad de Antioquia}\\
    \textit{Medellín-Colombia}

\end{center}

\vspace{0.5cm}

% -------------------- Resumen --------------------
\noindent
\textcolor{tituloazul}{\textbf{Resumen:}}
Esta práctica implementa un sistema de control por fase para la regulación de
potencia en una carga de corriente alterna (AC). El sistema emplea un
Rectificador Controlado de Silicio (SCR) BT151-500R como elemento principal
de conmutación, acompañado de un puente rectificador KBL04 que permite
controlar ambos semiciclos de la tensión de red mediante un único SCR.

La sincronización del sistema se obtiene mediante una etapa de detección de
cruce por cero, conformada por un transformador, un puente rectificador y un
optoacoplador RH111. La señal aislada obtenida es utilizada por el
microcontrolador para determinar el instante de disparo del SCR en función del
ángulo de fase $\alpha$. Posteriormente, el pulso de control activa un transistor
2N2222 y un optotriac MOC3020, proporcionando aislamiento galvánico entre la
etapa de control y la etapa de potencia.

Finalmente, se emplea un optoacoplador lineal IL300 junto con amplificadores
operacionales LM358 para realizar la medición aislada de la tensión sobre la
carga. De esta manera, se analiza experimentalmente la relación entre el
ángulo de disparo, el intervalo de conducción del SCR y la potencia entregada
a la carga.

% -------------------- Objetivo General --------------------
\section{Objetivo General}

Implementar un sistema de control de potencia por fase gestionado por un
microcontrolador y sincronizado mediante detección de cruce por cero,
utilizando un SCR para la regulación de energía en una carga AC.

% -------------------- Objetivos Específicos --------------------
\section{Objetivos Específicos}

\begin{itemize}
    \item Validar la etapa de detección de cruce por cero, garantizando la
    correcta sincronización de los pulsos del microcontrolador con la
    frecuencia de la red eléctrica.

    \item Establecer doble aislamiento galvánico mediante un transformador de
    medida y un optotriac, protegiendo la etapa de control frente a la red de
    potencia.

    \item Evaluar el comportamiento del SCR como interruptor mediante el
    análisis de la potencia entregada a la carga en función del ángulo de
    disparo $\alpha$.
\end{itemize}

% ============================================================
%                         PREINFORME
% ============================================================

\newpage

\section{Preinforme}

Se debe entregar/presentar lo siguiente al momento de entrar al laboratorio,
la falta de alguno de estos implicará una calificación de 0.0 en la práctica.

\subsection{Consulta de dispositivos}

\subsubsection{SCR BT151-500R}

El Rectificador Controlado de Silicio (SCR) es un dispositivo semiconductor
de tres terminales: ánodo (A), cátodo (K) y compuerta (G). Su funcionamiento
como interruptor se basa en el control de la corriente mediante la terminal
de compuerta.

En condición de bloqueo, el SCR puede soportar una tensión entre ánodo y
cátodo sin conducir una corriente significativa. Cuando se aplica un pulso
de corriente adecuado a la compuerta y el dispositivo se encuentra
polarizado directamente, el SCR entra en conducción.

Una característica fundamental es que, después del disparo, el SCR permanece
conduciendo aunque desaparezca el pulso de compuerta, siempre que la corriente
principal permanezca por encima de la corriente de mantenimiento $I_H$.

Por esta razón, para utilizarlo como interruptor se deben considerar
principalmente la tensión ánodo-cátodo, la corriente de carga y las
condiciones de disparo de compuerta.

\subsubsection{TRIAC MOC3020}

El TRIAC es un dispositivo semiconductor bidireccional que permite controlar
corriente en ambos sentidos. Debido a esta característica es ampliamente
utilizado en aplicaciones de control de potencia AC.

El MOC3020 es un optoacoplador cuya salida presenta comportamiento tipo
TRIAC. Su función en el circuito no es manejar directamente la potencia de
la carga, sino transferir el pulso de disparo desde la electrónica de control
hacia la etapa de potencia manteniendo aislamiento galvánico.

El MOC3020 es especialmente apropiado para control por fase debido a que es
un optotriac de disparo aleatorio. Esto permite que pueda activarse en un
instante determinado del semiciclo y no únicamente cerca del cruce por cero.

\subsubsection{Comparación entre SCR y TRIAC}

El SCR permite controlar la corriente principalmente en un solo sentido,
mientras que el TRIAC puede conducir corriente en ambos sentidos.

En una aplicación AC, un TRIAC puede utilizarse directamente para controlar
los dos semiciclos. En el circuito implementado, en cambio, se emplea un SCR
junto con un puente rectificador KBL04 para obtener el control de ambos
semiciclos mediante un dispositivo unidireccional.

\subsubsection{Optoacoplador lineal IL300}

El IL300 es un optoacoplador diseñado para la transferencia aislada de
señales analógicas. A diferencia de un optoacoplador digital convencional,
su estructura permite obtener una transferencia aproximadamente proporcional
entre la señal de entrada y la señal de salida.

El dispositivo incorpora un LED y fotodiodos que permiten realizar una
transferencia óptica con realimentación, reduciendo el efecto de variaciones
debidas a temperatura, envejecimiento y dispersión de las características
del LED.

En esta práctica, el IL300 permite obtener una representación aislada de la
tensión aplicada a la carga y enviarla posteriormente al microcontrolador
mediante una etapa de acondicionamiento basada en LM358.

\subsection{Código de control de fase}

El microcontrolador recibe una señal de sincronización proveniente de la
etapa de detección de cruce por cero. Cada vez que ocurre un cruce por cero,
se reinicia la temporización correspondiente al semiciclo y se espera un
intervalo determinado antes de generar el pulso de disparo.

El tiempo de retardo se relaciona con el ángulo de disparo mediante:

\begin{equation}
t_{\alpha} =
\frac{\alpha}{360^\circ}T
\label{eq:retardo}
\end{equation}

Para una frecuencia de red de \SI{60}{Hz}:

\begin{equation}
T=\frac{1}{f}
=\frac{1}{60}
\approx\SI{16.67}{ms}
\label{eq:periodo}
\end{equation}

Por ejemplo, para $\alpha=60^\circ$:

\begin{equation}
t_{\alpha}
=
\frac{60^\circ}{360^\circ}
\left(\SI{16.67}{ms}\right)
\approx\SI{2.78}{ms}
\label{eq:retardo60}
\end{equation}

Por lo tanto, el método utilizado corresponde técnicamente a un sistema de
control por fase sincronizado con la red y no a un PWM convencional de alta
frecuencia.

\subsection{Simulación}

El circuito de control de fase debe simularse verificando principalmente:

\begin{itemize}
    \item La detección del cruce por cero.
    \item El retardo correspondiente al ángulo de disparo.
    \item El funcionamiento del MOC3020.
    \item El disparo del SCR BT151-500R.
    \item La tensión obtenida sobre la carga.
    \item La tensión soportada por el SCR durante su estado de bloqueo.
\end{itemize}

La simulación debe permitir observar que el aumento del ángulo de disparo
reduce el intervalo de conducción del SCR y, por lo tanto, disminuye la
potencia entregada a la carga.

% ============================================================
%                         PRÁCTICA
% ============================================================

\section{Práctica}

Se implementó el circuito de control por fase siguiendo el protocolo de
seguridad establecido para la práctica. El sistema se dividió en las
etapas de detección de cruce por cero, generación del retardo, activación
aislada del SCR, etapa de potencia y medición aislada.

\subsection{Etapa de aislamiento y detección de cruce por cero}

El transformador recibe la tensión AC de la red mediante su devanado primario
y entrega una tensión reducida en el secundario identificado mediante los
terminales S1 y S2.

Esta etapa cumple tres funciones principales:

\begin{itemize}
    \item Reducir la tensión proveniente de la red a un nivel adecuado para
    el circuito de detección.

    \item Proporcionar aislamiento galvánico entre la red y la electrónica
    de control.

    \item Entregar una señal sincronizada con la red que permita detectar
    los cruces por cero.
\end{itemize}

Los diodos D1, D2, D3 y D4, implementados con 1N4007, forman un puente
rectificador. Esto permite que la corriente aplicada al LED interno del
RH111 mantenga la polaridad adecuada durante ambos semiciclos.

La resistencia R1 de \SI{1}{k\ohm} limita la corriente de la etapa de entrada,
mientras que el RH111 transfiere la información de la señal mediante
aislamiento óptico.

Cuando la tensión instantánea del secundario se aproxima a cero, la corriente
del LED interno disminuye y la salida del optoacoplador cambia de estado.
Esta transición permite obtener la señal digital ISR utilizada para
sincronizar el microcontrolador.

La resistencia R2 de \SI{470}{\ohm} establece la polarización de la salida,
mientras que R3 y D5 acondicionan y protegen la señal que llega a la entrada
del microcontrolador.

\subsection{Cálculo del ángulo de disparo}

El microcontrolador utiliza la señal ISR como referencia temporal. Al detectar
cada cruce por cero:

\begin{enumerate}
    \item Se reinicia la temporización del semiciclo.
    \item Se calcula el retardo correspondiente al ángulo $\alpha$.
    \item Se espera el tiempo calculado.
    \item Se genera el pulso denominado Disparo.
    \item El proceso se repite en el siguiente semiciclo.
\end{enumerate}

De acuerdo con \cref{eq:retardo}, aumentar $\alpha$ implica aumentar el tiempo
entre el cruce por cero y el pulso de Gate.

\subsection{Etapa de activación del MOC3020}

La salida del microcontrolador no activa directamente el SCR. La señal
Disparo pasa por una etapa formada por R4, R7, Q1, R14 y el MOC3020.

R4, de \SI{4.7}{k\ohm}, limita la corriente que llega a la base del transistor
Q1. R7, de \SI{10}{k\ohm}, mantiene la base del transistor en un estado
definido cuando la salida del microcontrolador se encuentra inactiva.

El transistor Q1, un 2N2222, funciona como interruptor de corriente. Cuando
la señal Disparo se encuentra en nivel alto, Q1 conduce y permite la
circulación de corriente por el LED interno del MOC3020.

La resistencia R14, de \SI{220}{\ohm}, limita la corriente del LED del
optoacoplador.

El MOC3020 proporciona aislamiento galvánico entre el circuito de control y
la etapa de potencia. Su salida tipo TRIAC permite transferir el pulso de
disparo en el instante requerido, característica necesaria para realizar
control de fase.

\subsection{Etapa de potencia}

El puente rectificador KBL04 permite utilizar un único SCR para controlar
ambos semiciclos de la red AC.

Durante cada semiciclo, un par de diodos del puente conduce y el sentido de
la corriente a través del SCR permanece constante. De esta manera, el SCR
BT151-500R puede utilizarse como interruptor unidireccional dentro de un
sistema de alimentación AC.

El comportamiento del SCR puede resumirse en tres estados:

\begin{enumerate}
    \item Inicialmente permanece bloqueado.
    \item Al recibir una corriente suficiente en Gate, se dispara.
    \item Una vez disparado, permanece conduciendo hasta que la corriente
    cae por debajo de la corriente de mantenimiento.
\end{enumerate}

Para una carga resistiva, el apagado ocurre aproximadamente cuando la
corriente alcanza nuevamente el cruce por cero.

\subsection{Medición aislada mediante IL300}

La tensión de la carga es acondicionada mediante la red formada por R11 y
R12 antes de ser aplicada a la etapa del IL300.

El LM358 ubicado en el lado de potencia acondiciona la señal y controla la
corriente del LED del IL300 mediante R8.

El IL300 transfiere la información de forma óptica y mantiene el aislamiento
galvánico entre la etapa de potencia y la etapa de control.

Posteriormente, otro LM358 acondiciona la señal obtenida en el lado aislado
para llevarla a la entrada analógica A0 del microcontrolador.

La trayectoria de medición puede resumirse como:

\[
V_{\mathrm{carga}}
\rightarrow
\text{divisor resistivo}
\rightarrow
\text{IL300}
\rightarrow
\text{LM358}
\rightarrow
A0.
\]

% ============================================================
%                  RESULTADOS Y ANÁLISIS
% ============================================================

\section{Informe (Resultados y Análisis)}

\subsection{Señal de control a la salida del microcontrolador y tensión
en la compuerta del SCR}

La salida del microcontrolador genera un pulso de control con un nivel
nominal de aproximadamente \SI{5}{V}. La posición temporal del pulso depende
del ángulo de disparo $\alpha$ establecido en el programa.

La señal atraviesa la siguiente cadena de activación:

\[
\text{Microcontrolador}
\rightarrow R4
\rightarrow Q1
\rightarrow MOC3020
\rightarrow R5
\rightarrow \text{Gate del SCR}.
\]

El transistor Q1 funciona como una etapa de excitación de corriente para el
LED interno del MOC3020. Esto evita exigir directamente al pin del
microcontrolador la corriente necesaria para activar el optotriac y mejora
la confiabilidad del disparo.

El MOC3020 transfiere ópticamente la orden hacia la etapa de potencia,
manteniendo el aislamiento galvánico. Por tanto, la señal del
microcontrolador representa la orden de disparo, mientras que la señal en
Gate corresponde a la activación final del SCR.

\subsection{Forma de onda de tensión sobre la carga y sobre el SCR}

Para analizar el comportamiento del sistema se considera un ángulo de
disparo de:

\[
\alpha=60^\circ.
\]

La tensión de alimentación utilizada es aproximadamente
\SI{120}{V RMS} a \SI{60}{Hz}. El valor pico correspondiente es:

\begin{equation}
V_{\mathrm{pico}}
=
V_{\mathrm{RMS}}\sqrt{2}
=
120\sqrt{2}
\approx
\SI{169.7}{V}.
\end{equation}

Por lo tanto, puede aproximarse a:

\[
V_{\mathrm{pico}}\approx\SI{170}{V}.
\]

El período de la red es:

\[
T=\frac{1}{60}
\approx\SI{16.67}{ms}.
\]

Para $\alpha=60^\circ$, el retardo de disparo es:

\[
t_\alpha
=
\frac{60^\circ}{360^\circ}
\left(\SI{16.67}{ms}\right)
\approx
\SI{2.78}{ms}.
\]

Durante el intervalo comprendido entre el cruce por cero y el instante
$t_\alpha$, el SCR permanece bloqueado y la tensión aplicada a la carga es
prácticamente nula.

Una vez aplicado el pulso de Gate, el SCR entra en conducción y la tensión
de la fuente aparece sobre la carga durante el resto del semiciclo.

De forma idealizada:

\begin{equation}
V_s(t)=V_m\sin(\omega t),
\end{equation}

donde:

\[
V_m\approx\SI{170}{V}
\]

y

\[
\omega=2\pi f.
\]

Por tanto, el control de fase produce un recorte de la señal de tensión
aplicada a la carga.

Al aumentar $\alpha$, disminuye el intervalo de conducción:

\[
\alpha\uparrow
\quad\Rightarrow\quad
V_{\mathrm{carga,RMS}}\downarrow
\]

mientras que:

\[
\alpha\downarrow
\quad\Rightarrow\quad
V_{\mathrm{carga,RMS}}\uparrow.
\]

La tensión sobre el SCR presenta un comportamiento complementario. Mientras
el SCR está bloqueado soporta una parte importante de la tensión de la
etapa de potencia. Una vez que entra en conducción, la tensión ánodo-cátodo
disminuye considerablemente y queda aproximadamente en el rango de la
caída de tensión propia del dispositivo.

Para las condiciones de operación consideradas puede tomarse como
referencia:

\[
V_{AK(\mathrm{on})}\approx\SI{1}{V}-\SI{2}{V}.
\]

\subsection{Comparación entre la tensión de carga y la señal del IL300}

El IL300 tiene una función diferente a la del MOC3020.

El MOC3020 transmite la orden de disparo desde la etapa de control hacia la
etapa de potencia, mientras que el IL300 permite transmitir información
analógica en sentido contrario.

La trayectoria de medición es:

\[
V_{\mathrm{carga}}
\rightarrow
R11-R12
\rightarrow
IL300
\rightarrow
LM358
\rightarrow
A0.
\]

La red resistiva reduce la tensión proveniente de la carga hasta un nivel
adecuado para el circuito de medición. Posteriormente, el IL300 realiza la
transferencia óptica manteniendo el aislamiento galvánico.

El LM358 acondiciona la señal para adaptarla al rango de entrada analógica
del microcontrolador.

El valor pico mostrado por la interfaz serial depende de la tensión real de
la carga, de la relación de reducción de R11-R12, de la transferencia del
IL300 y de la ganancia configurada mediante el LM358.

\subsection{Diferencia entre el IL300 y el MOC3020}

\begin{table}[H]
\centering
\caption{Comparación entre el MOC3020 y el IL300.}
\label{tComparacionOptos}
\begin{tabular}{p{4cm} p{5cm} p{5cm}}
\toprule
\textit{Característica} & \textit{MOC3020} & \textit{IL300} \\
\midrule
Tipo &
Optoacoplador con salida tipo TRIAC &
Optoacoplador lineal \\

Función &
Disparo &
Medición y realimentación \\

Tipo de señal &
Pulso de control &
Señal analógica \\

Dirección &
Control $\rightarrow$ potencia &
Potencia $\rightarrow$ control \\

Aplicación &
Activación del circuito de disparo &
Medición aislada de la tensión de carga \\

Aislamiento galvánico &
Sí &
Sí \\

\bottomrule
\end{tabular}
\end{table}

En síntesis, el MOC3020 se utiliza para enviar la orden de disparo,
mientras que el IL300 permite informar al microcontrolador sobre el
comportamiento de la señal de potencia.

\subsection{Variación del ángulo de disparo}

Al modificar el ángulo de disparo desde el microcontrolador se modifica el
instante en que el SCR comienza a conducir.

Para un ángulo pequeño, el disparo ocurre poco después del cruce por cero,
por lo que el SCR permanece conduciendo durante una mayor parte del
semiciclo. Como consecuencia, se entrega una mayor cantidad de energía a
la carga.

Cuando el ángulo aumenta, el disparo ocurre posteriormente y disminuye el
tiempo de conducción. Por tanto, disminuyen la tensión RMS y la potencia
entregada.

El comportamiento general es:

\[
\alpha\downarrow
\quad\Rightarrow\quad
\text{mayor tiempo de conducción}
\]

\[
\alpha\uparrow
\quad\Rightarrow\quad
\text{menor tiempo de conducción}.
\]

Para una carga resistiva y una alimentación de \SI{120}{V RMS}, se tienen
los siguientes valores teóricos de referencia:

\begin{table}[H]
\centering
\caption{Tensión RMS teórica para diferentes ángulos de disparo.}
\label{tRMS}
\begin{tabular}{cc}
\toprule
\textit{Ángulo de disparo $\alpha$} &
\textit{Tensión RMS aproximada} \\
\midrule
$30^\circ$ & \SI{118.5}{V} \\
$60^\circ$ & \SI{107.8}{V} \\
$90^\circ$ & \SI{85.0}{V} \\
$120^\circ$ & \SI{53.2}{V} \\
\bottomrule
\end{tabular}
\end{table}

Estos valores corresponden a una referencia teórica y no deben interpretarse
como mediciones experimentales.

\subsection{Efecto de un pulso demasiado corto en el MOC3020}

El disparo confiable del SCR requiere una corriente de compuerta suficiente
durante un tiempo adecuado.

La condición básica de disparo puede expresarse como:

\begin{equation}
I_G\geq I_{GT}.
\end{equation}

Si el pulso enviado mediante el MOC3020 es demasiado corto, pueden
presentarse diferentes comportamientos:

\begin{itemize}
    \item El SCR puede no dispararse.
    \item El disparo puede ser intermitente.
    \item El SCR puede activarse en algunos ciclos y no en otros.
    \item La forma de onda de la carga puede presentar irregularidades.
    \item La potencia entregada puede variar entre ciclos.
\end{itemize}

Además, para que el SCR quede enclavado después del disparo, la corriente
principal debe alcanzar la corriente de enclavamiento $I_L$. Una vez
establecida la conducción, el SCR puede permanecer activado mientras la
corriente sea superior a la corriente de mantenimiento $I_H$.

Por esta razón, el diseño del pulso debe considerar simultáneamente la
corriente de Gate y su duración.

\subsection{Necesidad de implementar tres aislamientos galvánicos}

El sistema emplea tres mecanismos de aislamiento con funciones diferentes.

\begin{enumerate}
    \item \textbf{Transformador T1:} proporciona aislamiento entre la red
    eléctrica y la etapa de detección de sincronismo utilizada por el
    microcontrolador.

    \item \textbf{MOC3020:} proporciona aislamiento entre el pulso de control
    del microcontrolador y el circuito de disparo del SCR.

    \item \textbf{IL300:} proporciona aislamiento entre la tensión de la
    etapa de potencia y la entrada analógica del microcontrolador.
\end{enumerate}

Los tres aislamientos permiten transmitir información entre las diferentes
etapas sin establecer una conexión eléctrica directa.

Esto mejora la seguridad del sistema y reduce la posibilidad de que
sobretensiones, transitorios o corrientes de retorno provenientes de la
etapa de potencia alcancen el microcontrolador, el computador o los equipos
de medición.

\subsection{Separación entre la tierra de control y la referencia de potencia}

La tierra de control y la referencia de potencia deben permanecer
eléctricamente separadas.

La etapa de potencia trabaja con una tensión de aproximadamente
\SI{120}{V RMS}, cuyo valor pico es cercano a:

\[
V_{\mathrm{pico}}\approx\SI{170}{V}.
\]

En contraste, la electrónica de control trabaja con niveles del orden de
\SI{5}{V}.

Si ambas referencias fueran conectadas directamente, se perdería el
aislamiento galvánico y cualquier falla, sobretensión o transitorio de la
etapa de potencia podría propagarse hacia el microcontrolador.

Esto podría producir:

\begin{itemize}
    \item Daños en el microcontrolador.
    \item Corrientes de retorno no deseadas.
    \item Errores de medición.
    \item Cortocircuitos.
    \item Riesgo eléctrico para el usuario y los equipos.
\end{itemize}

Por tanto:

\begin{equation}
GND_{\mathrm{control}}
\neq
\text{Referencia}_{\mathrm{potencia}}.
\end{equation}

Los transformadores y optoacopladores permiten intercambiar información
sin eliminar esta separación eléctrica.

\subsection{Funcionamiento del optoacoplador lineal IL300}

El IL300 permite transferir una señal analógica manteniendo aislamiento
galvánico.

La tensión de la carga primero es reducida mediante el divisor resistivo
formado por R11 y R12. Esta señal acondicionada controla el LED interno del
IL300.

La corriente del LED produce una señal óptica que es detectada por los
fotodiodos internos del dispositivo. La estructura de realimentación del
IL300 permite obtener una transferencia más precisa que la de un
optoacoplador convencional utilizado como interruptor.

Posteriormente, el LM358 acondiciona la señal obtenida para adaptarla a la
entrada analógica A0 del microcontrolador.

La cadena completa es:

\[
V_{\mathrm{carga}}
\rightarrow
\text{divisor resistivo}
\rightarrow
IL300
\rightarrow
LM358
\rightarrow
A0.
\]

De esta forma, el microcontrolador puede obtener información relacionada
con la forma de onda de la carga sin estar conectado eléctricamente a la
etapa de potencia.

\subsection{Diferencia operativa entre SCR y TRIAC}

El SCR es un dispositivo controlado que conduce principalmente en un solo
sentido. Después de recibir un pulso adecuado en la compuerta, permanece
conduciendo mientras la corriente principal sea suficiente para mantener
su estado de conducción.

El TRIAC puede conducir corriente en ambos sentidos y, por ello, resulta
apropiado para controlar directamente cargas AC.

En el circuito estudiado, el MOC3020 posee una salida tipo TRIAC, pero no
maneja directamente la potencia de la carga. Su función es transferir de
forma aislada el pulso que participa en el disparo del SCR.

Por tanto:

\[
\boxed{\text{SCR: conducción principalmente unidireccional}}
\]

\[
\boxed{\text{TRIAC: conducción bidireccional}}
\]

El puente KBL04 permite compensar la característica unidireccional del SCR
y posibilita el control de ambos semiciclos de la alimentación AC.

\subsection{Análisis general de los resultados}

Los resultados obtenidos permiten verificar el principio de funcionamiento
del control de fase mediante el SCR BT151-500R.

La variación del ángulo de disparo modifica directamente el intervalo de
conducción del tiristor. Cuando $\alpha$ disminuye, el pulso de Gate aparece
más cerca del cruce por cero, permitiendo que el SCR conduzca durante una
mayor parte del semiciclo. En consecuencia, aumenta la tensión RMS aplicada
a la carga y, para una carga resistiva, también aumenta la potencia
transferida.

Por el contrario, un incremento de $\alpha$ desplaza el disparo hacia una
posición posterior dentro del semiciclo. Esto reduce el intervalo de
conducción y disminuye la tensión RMS y la potencia entregada.

El MOC3020 permite realizar este disparo manteniendo el aislamiento
galvánico entre el microcontrolador y la etapa de potencia. El transistor
Q1 proporciona la corriente necesaria para activar el LED interno del
MOC3020, evitando que el microcontrolador tenga que suministrarla
directamente.

El IL300 cumple una función complementaria, ya que permite obtener una
realimentación analógica aislada de la tensión presente en la carga. Esta
señal es acondicionada mediante LM358 y llevada a la entrada analógica del
microcontrolador.

Finalmente, la duración del pulso de Gate es un parámetro importante para
garantizar un disparo estable del SCR. Si la corriente de compuerta es
insuficiente o el pulso tiene una duración demasiado corta, puede producirse
un disparo intermitente y, por tanto, una potencia irregular en la carga.

En conjunto, el circuito integra sincronización con la red, control de
ángulo de fase, aislamiento galvánico y realimentación analógica, permitiendo
regular de manera controlada la potencia entregada a una carga AC.

% ============================================================
%                         CONCLUSIONES
% ============================================================

\section{Conclusiones}

\begin{itemize}

    \item Se verificó el principio de control de potencia por ángulo de fase
    utilizando el SCR BT151-500R. La posición temporal del pulso de Gate
    respecto al cruce por cero determina el instante de inicio de conducción
    del dispositivo y, por tanto, el intervalo durante el cual la carga
    recibe energía.

    \item La etapa de detección de cruce por cero permite establecer una
    referencia temporal sincronizada con la red eléctrica. Esta sincronización
    es fundamental para que el microcontrolador pueda generar retardos
    correspondientes a diferentes valores del ángulo de disparo $\alpha$.

    \item El MOC3020 permite transferir el pulso de disparo manteniendo
    aislamiento galvánico entre el sistema de control y la etapa de potencia.
    El transistor Q1 actúa como etapa de excitación, proporcionando la
    corriente necesaria para activar el LED interno del optotriac.

    \item El puente rectificador KBL04 permite emplear un único SCR para
    controlar ambos semiciclos de la señal AC. De esta manera, la corriente
    que atraviesa el SCR mantiene la misma dirección aunque la polaridad de
    la red cambie.

    \item Se estableció que el aumento del ángulo de disparo reduce el
    intervalo de conducción del SCR. Para una carga resistiva, esta reducción
    implica una disminución de la tensión RMS y de la potencia entregada a
    la carga.

    \item El IL300 permite realizar una medición analógica aislada de la
    tensión de carga. Su utilización junto con los amplificadores LM358
    permite acondicionar la señal y llevarla al microcontrolador sin
    establecer una conexión eléctrica directa entre la etapa de potencia y
    la electrónica de control.

    \item La implementación de los tres mecanismos de aislamiento,
    correspondientes al transformador, al MOC3020 y al IL300, permite
    separar las funciones de sincronización, disparo y medición, reduciendo
    el riesgo de transferencia de sobretensiones y transitorios hacia el
    microcontrolador y los equipos de medición.

    \item Finalmente, se determinó que un pulso de Gate debe presentar una
    corriente y duración adecuadas para garantizar el disparo y enclavamiento
    del SCR. Un pulso demasiado corto puede producir disparos intermitentes,
    deformación de la señal de carga y variaciones de potencia entre
    diferentes ciclos.

\end{itemize}

% ============================================================
%                         BIBLIOGRAFÍA
% ============================================================

\begin{thebibliography}{99}

\bibitem{ltspice}
Analog Devices,
``LTspice Simulator,''
2026.
Disponible en:
\url{https://www.analog.com/en/resources/design-tools-and-calculators/ltspice-simulator.html}.

\end{thebibliography}

% ============================================================
%                  EQUIPOS Y COMPONENTES
% ============================================================

\newpage

\section*{Listado de Equipos, Herramientas y Componentes}

{\textbf{\estudiantes}}

\begin{table}[H]
    \centering
    \caption{Elementos utilizados o solicitados para la Práctica 2.}

    \resizebox{0.9\textwidth}{!}{
    \begin{tabular}{c p{11cm} c}
    \toprule
    \toprule
    \textit{Item} & \textit{Elemento} & \textit{Cantidad} \\
    \midrule

    1 & Osciloscopio + puntas de medición & 1 \\

    \midrule
    2 & Multímetro + puntas de medición & 1 \\

    \midrule
    3 & Fuente DC & 1 \\

    \midrule
    4 & Protoboard & 1 \\

    \midrule
    5 & Transformador de aislamiento/medición & 1 \\

    \midrule
    6 & Microcontrolador (STM32, ESP32 o Raspberry Pi Pico) & 1 \\

    \midrule
    7 & SCR BT151-500R & 1 \\

    \midrule
    8 & Optotriac MOC3020 & 1 \\

    \midrule
    9 & Optoacoplador lineal IL300 & 1 \\

    \midrule
    10 & Optoacoplador RH111 & 1 \\

    \midrule
    11 & Puente rectificador KBL04 & 1 \\

    \midrule
    12 & Amplificador operacional LM358 & 2 \\

    \midrule
    13 & Transistor 2N2222 & 1 \\

    \midrule
    14 & Diodo 1N4007 & 4 \\

    \midrule
    15 & Diodo 1N4148 & 1 \\

    \midrule
    16 & Resistencia \SI{1}{k\ohm} & 2 \\

    \midrule
    17 & Resistencia \SI{470}{\ohm} & 1 \\

    \midrule
    18 & Resistencia \SI{4.7}{k\ohm} & 2 \\

    \midrule
    19 & Resistencia \SI{10}{k\ohm} & 1 \\

    \midrule
    20 & Resistencia \SI{220}{\ohm} & 2 \\

    \midrule
    21 & Bombillo AC & 1 \\

    \midrule
    22 & Alambre para conexiones & Según necesidad \\

    \midrule

    \bottomrule
    \end{tabular}
    }

    \label{tElementos}

\end{table}

\end{document}
